\documentclass[12pt]{article}
\usepackage[T1]{fontenc}
\usepackage[utf8]{inputenc}
\usepackage{lmodern}
\usepackage{textcomp}
\usepackage{enumitem}
\usepackage{geometry}
\geometry{margin=1in}
\begin{document}
\begin{center}
Answer Key - Geography - K-12 Level Assessment 10 \
\small (Answers are ordered exactly as in the question paper.)
\end{center}

\section*{Multiple Choice}
\begin{enumerate}[label=\arabic*.]
\item \textbf{Answer:} B
\\ \textit{Options:} A) the change of government from a democracy to a communist dictatorship.; B) the movement of power from the central government to regional governments in the country.; C) the movement of power from a monarchy to a military dictatorship.; D) the movement of power from regional governments to central governments.
\\ \textit{Note:} The correct answer describes devolution as the transfer of authority and decision-making power from a central government to regional or local governments, reflecting the key concept of decentralization in political geography.
\item \textbf{Answer:} B
\\ \textit{Options:} A) Truck; B) Air; C) Ship; D) Railroad
\\ \textit{Note:} Shipping the jewelry by air is the most sensible choice because it offers the fastest delivery time and enhanced security, which is crucial for transporting valuable items.
\item \textbf{Answer:} C
\\ \textit{Options:} A) central business districts and suburbs; B) a market and its sources for raw materials; C) cities of different sizes and functions; D) the hubs of airline transportation systems
\\ \textit{Note:} Walter Christaller's central place theory analyzes how cities of varying sizes and functions are strategically located within a region to serve as central places for surrounding areas, optimizing accessibility and economic efficiency.
\item \textbf{Answer:} A
\\ \textit{Options:} A) Rice; B) Wheat; C) Barley; D) Millet
\\ \textit{Note:} Rice is considered the crop that began the Third Agricultural Revolution because its domestication and the development of high-yield varieties significantly increased food production and agricultural efficiency, paving the way for advancements in farming practices and technology.
\item \textbf{Answer:} D
\\ \textit{Options:} A) An industrial society with abundant natural resources and large imports of food; B) A society with a highly mechanized agricultural sector; C) A non-ecumene; D) A slash-and-burn agricultural society
\\ \textit{Note:} A slash-and-burn agricultural society is most likely to experience population pressure because its farming method relies on clearing land for cultivation, which can quickly degrade soil fertility and limit sustainable agricultural practices as the population grows.
\end{enumerate}

\section*{True / False}
\begin{enumerate}[label=\arabic*.]
\setcounter{enumi}{5}
\item \textbf{Answer:} False
\\ \textit{Options:} A) True; B) False
\\ \textit{Note:} Deglomeration refers to the process of dispersing production activities and people away from concentrated areas, often to reduce costs or alleviate congestion, rather than bringing them closer together for economic growth.
\item \textbf{Answer:} True
\\ \textit{Options:} A) True; B) False
\\ \textit{Note:} Palestine is considered a stateless nation because it lacks full sovereignty and control over its territory, as its governance is limited by ongoing conflict and occupation.
\item \textbf{Answer:} False
\\ \textit{Options:} A) True; B) False
\\ \textit{Note:} Farm workers often engage in circular migration, where they temporarily move to work in host countries rather than permanently settling there with their families.
\item \textbf{Answer:} True
\\ \textit{Options:} A) True; B) False
\\ \textit{Note:} Quebec is considered a prominent example of ethnonationalism because it has a unique French-speaking cultural identity that seeks greater autonomy and recognition within the multicultural framework of Canada.
\item \textbf{Answer:} True
\\ \textit{Options:} A) True; B) False
\\ \textit{Note:} NATO is primarily focused on collective defense and military cooperation among member countries, whereas the EU and NAFTA emphasize economic integration and trade agreements.
\end{enumerate}

\section*{Long Form Response}
\begin{enumerate}[label=\arabic*.]
\setcounter{enumi}{10}
\item \textbf{Answer:} Replacement level fertility refers to the number of children a couple must have to replace themselves in the population, typically estimated at about 2.1 children per woman in developed countries. This concept is significant in population dynamics because it helps determine whether a population is growing, stable, or declining. When fertility rates fall below this level, it can lead to an aging population and potential labor shortages, impacting economic sustainability and social services. Conversely, if fertility rates are above replacement level, it can lead to rapid population growth, which may strain resources and the environment. Understanding replacement level fertility is crucial for policymakers to address demographic trends and ensure a balanced and sustainable society.
\\ \textit{Note:} correct because it accurately defines replacement level fertility, explains its implications for population growth or decline, and highlights its significance in influencing economic and social sustainability within a society.
\item \textbf{Answer:} The end of the Cold War in the early 1990 s marked a significant shift in political dynamics, leading to the destabilization of Eastern Europe. With the collapse of the Soviet Union, many countries in the region sought independence, which sparked nationalistic feelings among various ethnic groups. As borders were redrawn and new nations emerged, long-standing ethnic tensions resurfaced, resulting in conflicts over territory and power. This fragmentation often led to violence, as different groups vied for control and recognition. Consequently, the end of the Cold War not only changed political alliances but also ignited complex ethnic conflicts that have shaped Eastern Europe over the past 25 years.
\\ \textit{Note:} The answer correctly identifies that the end of the Cold War led to the collapse of the Soviet Union, which triggered nationalistic movements and the resurgence of ethnic tensions in Eastern Europe, resulting in conflicts over territory and power among newly independent nations.
\end{enumerate}

\vfill
\noindent\textit{End of Answer Key}
\end{document}